	%	--------------------------------------------------------------
% thesis.tex 
%--------------------------------------------------------------
% Corso di Laurea in Informatica 
% http://if.dsi.unifi.it/
% @Facolt\`a di Scienze Matematiche, Fisiche e Naturali
% @Universit\`a degli Studi di Firenze
%--------------------------------------------------------------
% - template for the main file of Informatica@Unifi Thesis 
% - based on Classic Thesis Style Copyright (C) 2008 
%   Andr\'e Miede http://www.miede.de   
%--------------------------------------------------------------
\documentclass[twoside,openright,titlepage,fleqn,
	headinclude,12pt,a4paper,BCOR5mm,footinclude]{scrbook}
%--------------------------------------------------------------
\newcommand{\myItalianTitle}{Composizione di operatori di ordinamento con contenitori\xspace}
\newcommand{\myEnglishTitle}{Composition of sorting operators with containers\xspace}
% use the right myDegree option
\newcommand{\myDegree}{Corso di Laurea in Informatica\xspace}
%\newcommand{\myDegree}{
	%Corso di Laurea Specialistica in Scienze e Tecnologie 
	%dell'Informazione\xspace}
\newcommand{\myName}{Alessio Santoro\xspace}
\newcommand{\myProf}{Relatore\xspace}
\newcommand{\myOtherProf}{Correlatore\xspace}
\newcommand{\mySupervisor}{Nome Cognome\xspace}
\newcommand{\myFaculty}{
	Scuola di Scienze Matematiche, Fisiche e Naturali\xspace}
\newcommand{\myUni}{\protect{
	Universit\`a degli Studi di Firenze}\xspace}
\newcommand{\myLocation}{Firenze\xspace}
\newcommand{\myTime}{Anno Accademico 2023-2024\xspace}
\newcommand{\myVersion}{Version 0.1\xspace}
%--------------------------------------------------------------
\usepackage[italian]{babel}
\usepackage[latin1]{inputenc} 
\usepackage[T1]{fontenc} 
\usepackage[square,numbers]{natbib} 
\usepackage[fleqn]{amsmath}  
\usepackage{ellipsis}
\usepackage{listings}
\usepackage{subfig}
\usepackage{caption}
\usepackage{appendix}
\usepackage{siunitx}
\usepackage{amsthm}
\usepackage{algorithm}
\usepackage{algpseudocode}
%--------------------------------------------------------------
\usepackage{dia-classicthesis-ldpkg}
%\usepackage{classicthesis}
%--------------------------------------------------------------
% Options for classicthesis.sty:
% tocaligned eulerchapternumbers drafting linedheaders 
% listsseparated subfig nochapters beramono eulermath parts 
% minionpro pdfspacing
\usepackage[eulerchapternumbers,linedheaders,subfig,beramono,eulermath,
parts]{classicthesis}
%--------------------------------------------------------------
\newlength{\abcd} % for ab..z string length calculation
% how all the floats will be aligned
\newcommand{\myfloatalign}{\centering} 
\setlength{\extrarowheight}{3pt} % increase table row height
\captionsetup{format=hang,font=small}
%--------------------------------------------------------------
% Layout setting
%--------------------------------------------------------------
\usepackage{geometry}
\geometry{
	a4paper,
	ignoremp,
	bindingoffset = 1cm, 
	textwidth     = 13.5cm,
	textheight    = 21.5cm,
	lmargin       = 3.5cm, % left margin
	tmargin       = 4cm    % top margin 
}

\lstset{
  	frame=tb,
	language=Matlab,
  	aboveskip=3mm,
  	belowskip=3mm,
  	showstringspaces=false,
  	columns=flexible,
  	basicstyle={\small\ttfamily},
  	numbers=none,
  	breaklines=true,
  	breakatwhitespace=true,
  	tabsize=3
}
\usepackage{tikz}
\usetikzlibrary{matrix,trees,arrows}

\usepackage{stmaryrd}		% various mathematical symbols

% For tikz pictures
\usetikzlibrary{positioning}
\usetikzlibrary{fit}
\usetikzlibrary{patterns}

\newcommand{\ZZ}{\mathbb{Z}}           % Integers
\newcommand{\sym}{\mathcal{S}}         % Permutations
\DeclareMathOperator{\unSvar}{{\mathsf{cand}}}
\DeclareMathOperator{\inv}{{\mathsf{inv}}}
\DeclareMathOperator{\ninv}{{\mathsf{ninv}}}
\DeclareMathOperator{\symS}{{\mathcal{S}}}

\newcommand{\dbrac}[1]{{\llbracket #1 \rrbracket}} 	% double brackets

\DeclareMathOperator{\Av}{\mathrm{Av}}

%%%%%%%%%%% MACROS FOR DRAWING INTERVAL AND MESH PATTERNS %%%%%%%%%%%

% Sub-macros
\newcommand{\shadetheboxes}[1]{
	\foreach \x/\y in {#1}
      	\fill[pattern color = black!65, pattern=north east lines] (\x,\y) rectangle +(1,1);
	}
	
\newcommand{\drawthegrid}[1]{
	\draw (0.01,0.01) grid (#1+0.99,#1+0.99);
	}
	
\newcommand{\drawtheclpattern}[1]{
	\foreach \x/\y in {#1}
      	\filldraw (\x,\y) circle (6pt);
	}
	
\newcommand{\drawtheclpatternwhite}[1]{
	\foreach \x/\y in {#1}
      	\draw[fill=white] (\x,\y) circle (6pt);
	}
	
\newcommand{\drawtheclpatternwhitebig}[1]{
	\foreach \x/\y in {#1}
		\draw[fill=white] (\x,\y) circle (11pt);
	}
	
\newcommand{\drawspecialbox}[1]{
	\foreach \x/\y/\z/\w/\A in {#1}
		{
       		\fill[color = white!100, opacity=1, rounded corners = 1.5pt] (\x+0.125,\y+0.125) rectangle (\z-0.125,\w-0.125);
       		\draw[color = black, rounded corners = 1.5pt] (\x+0.125,\y+0.125) rectangle (\z-0.125,\w-0.125);
       		\fill[black] (\x/2+\z/2,\y/2+\w/2) node {$\scriptstyle\A$};
       	}
    }
    
\newcommand{\drawspecialboxnomark}[1]{
	\foreach \x/\y/\z/\w in {#1}
		{
       		\fill[color = white!100, opacity=1, rounded corners = 1.5pt] (\x+0.125,\y+0.125) rectangle (\z-0.125,\w-0.125);
       		\draw[color = black, rounded corners = 1.5pt] (\x+0.125,\y+0.125) rectangle (\z-0.125,\w-0.125);
       	}
    }
    
\newcommand{\drawspecialboxthemark}[1]{
	\foreach \x/\y/\z/\w/\A in {#1}
		\fill[black] (\x/2+\z/2,\y/2+\w/2) node {$\scriptstyle\A$};
    }
    
\newcommand{\drawspecialshadedbox}[1]{
	\foreach \x/\y/\z/\w/\A in {#1}
		{
       		\fill[color = white!100, opacity=1, rounded corners=1.5pt] (\x+0.125,\y+0.125) rectangle (\z-0.125,\w-0.125);
       		\fill[pattern color = black!65, pattern=north east lines, rounded corners=1.5pt] (\x+0.125,\y+0.125) rectangle (\z-0.125,\w-0.125);
       		\draw[color = black, rounded corners=1.5pt] (\x+0.125,\y+0.125) rectangle (\z-0.125,\w-0.125);
       		\fill[black] (\x/2+\z/2,\y/2+\w/2) node {$\scriptstyle\A$};
       	}
    }
    

\newcommand{\mpattern}[4]{										% mesh pattern
  \raisebox{0.6ex}{
  \begin{tikzpicture}[scale=0.35, baseline=(current bounding box.center), #1]
  	\useasboundingbox (0.0,-0.1) rectangle (#2+1.4,#2+1.1);
	
    \shadetheboxes{#4}
    
    \drawthegrid{#2}
    
    \drawtheclpattern{#3}
    
  \end{tikzpicture}}
}

\newcommand{\mpatternnl}[4]{										% mesh pattern - nl
  \raisebox{0.6ex}{
  \begin{tikzpicture}[scale=0.35, baseline=(current bounding box.center), #1]
  	\useasboundingbox (0.85,-0.1) rectangle (#2+1.4,#2+1.1);
	
    \shadetheboxes{#4}
    
    \drawtheclpattern{#3}
    
  \end{tikzpicture}}
}

\newcommand{\onetwo}{\mpatternnl{scale=0.2}{2}{1/1,2/2}{}}		% the pattern 12 with no lines

\newcommand{\mpatternww}[5]{									% mesh pattern with white and black dots
  \raisebox{0.6ex}{
  \begin{tikzpicture}[scale=0.35, baseline=(current bounding box.center), #1]
  	\useasboundingbox (0.0,-0.1) rectangle (#2+1.4,#2+1.1);
	
    \shadetheboxes{#5}
    
    \drawthegrid{#2}
    
    \drawtheclpatternwhite{#4}
    \drawtheclpattern{#3}
    
  \end{tikzpicture}}
}

\newcommand{\mpatternwwo}[6]{									% mesh pattern with white and black dots overlapping
  \raisebox{0.6ex}{
  \begin{tikzpicture}[scale=0.35, baseline=(current bounding box.center), #1]
  \useasboundingbox (0.0,-0.1) rectangle (#2+1.4,#2+1.1);
  
    \shadetheboxes{#6}
    
    \drawthegrid{#2}
    
    \drawtheclpatternwhite{#4}
    \drawtheclpatternwhitebig{#5}
    \drawtheclpattern{#3}
    
  \end{tikzpicture}}
}

\newcommand{\mmpattern}[5]{									% mesh patterns with a special box marked
  \raisebox{0.6ex}{
  \begin{tikzpicture}[scale=0.35, baseline=(current bounding box.center), #1]
  \useasboundingbox (0.0,-0.1) rectangle (#2+1.4,#2+1.1);
    
    \shadetheboxes{#4}
    
    \drawthegrid{#2}
    
    \drawspecialbox{#5}
    
    \drawtheclpattern{#3}

  \end{tikzpicture}}
}

\newcommand{\mmpatternwwo}[7]{									% mesh pattern with white and black dots overlapping
  \raisebox{0.6ex}{
  \begin{tikzpicture}[scale=0.35, baseline=(current bounding box.center), #1]
  \useasboundingbox (0.0,-0.1) rectangle (#2+1.4,#2+1.1);
    
    \shadetheboxes{#6}
    
    \drawthegrid{#2}
    
    \drawspecialbox{#7}
    
    \drawtheclpatternwhite{#4}
    \drawtheclpatternwhitebig{#5}
    \drawtheclpattern{#3}
    
  \end{tikzpicture}}
}

\newcommand{\mmpatternloc}[6]{									% mesh patterns with a special box marked
  \raisebox{0.6ex}{											% can control location of mark
  \begin{tikzpicture}[scale=0.35, baseline=(current bounding box.center), #1]
  \useasboundingbox (0.0,-0.1) rectangle (#2+1.4,#2+1.1);
    
    \shadetheboxes{#4}
    
    \drawthegrid{#2}
    
    \drawspecialboxnomark{#5}
    \drawspecialboxthemark{#6}
    
    \drawtheclpattern{#3}

  \end{tikzpicture}}
}

\newcommand{\mmpatternwwoloc}[8]{									% mesh pattern with white and black dots overlapping
  \raisebox{0.6ex}{												% can control location of mark
  \begin{tikzpicture}[scale=0.35, baseline=(current bounding box.center), #1]
  \useasboundingbox (0.0,-0.1) rectangle (#2+1.4,#2+1.1);
    
    \shadetheboxes{#6}
    
    \drawthegrid{#2}
    
    \drawspecialboxnomark{#7}
    \drawspecialboxthemark{#8}
    
    \drawtheclpatternwhite{#4}
    \drawtheclpatternwhitebig{#5}
    \drawtheclpattern{#3}
    
  \end{tikzpicture}}
}

\newcommand{\decpattern}[6]{									% decorated pattern
  \raisebox{0.6ex}{
  \begin{tikzpicture}[scale=0.35, baseline=(current bounding box.center), #1]
  \useasboundingbox (0.0,-0.1) rectangle (#2+1.4,#2+1.1);
    
    \shadetheboxes{#4}
    
    \drawthegrid{#2}
       
    \drawspecialbox{#5}
    \drawspecialshadedbox{#6}
    
    \drawtheclpattern{#3}

  \end{tikzpicture}}
}

\newcommand{\decpatternww}[8]{								% decorated pattern w white
  \raisebox{0.6ex}{
  \begin{tikzpicture}[scale=0.35, baseline=(current bounding box.center), #1]
  \useasboundingbox (0.0,-0.1) rectangle (#2+1.4,#2+1.1);
    
    \shadetheboxes{#6}
    
    \drawthegrid{#2}
       
    \drawspecialbox{#7}
    \drawspecialshadedbox{#8}
    
    \drawtheclpatternwhite{#4}
    \drawtheclpatternwhitebig{#5}
    \drawtheclpattern{#3}

  \end{tikzpicture}}
}

\newcommand{\mpatternlabelled}[5]{										% mesh pattern with labels for some boxes
  \raisebox{0.6ex}{
  \begin{tikzpicture}[scale=0.35, baseline=(current bounding box.center), #1]
    
    \shadetheboxes{#4}
    
    \drawthegrid{#2}
    
    \drawtheclpattern{#3}
    
    \draw (0.5,0.5) -- (0.5,-0.5);
    \fill[black] (0.5,-1) node {$\scriptstyle(0,0)$};
    \draw (0.5,2.5) -- (0.5,4.5);
    \fill[black] (0.5,5) node {$\scriptstyle(0,2)$};
    \draw (1.5,2.5) -- (1.5,6);
    \fill[black] (1.5,6.5) node {$\scriptstyle(1,2)$};
    \draw (2.5,2.5) -- (2.5,4.5);
    \fill[black] (2.5,5) node {$\scriptstyle(2,2)$};
    
  \end{tikzpicture}}
}

%--------------------------------------------------------------
\begin{document}
\frenchspacing
\raggedbottom
\pagenumbering{roman}
\pagestyle{plain}
%--------------------------------------------------------------
% Frontmatter
%--------------------------------------------------------------
%--------------------------------------------------------------
% titlepage.tex (use thesis.tex as main file)
%--------------------------------------------------------------
\begin{titlepage}
	\begin{center}
   	\large
      \hfill
      \vfill
      \begingroup
         \includegraphics[scale=0.15]{logo/LOGO}\\
%			\spacedallcaps{\myUni} \\ 
			\myFaculty \\
			\myDegree \\ 
			\vspace{0.5cm}
         \vspace{0.5cm}    
         Tesi di Laurea    
      \endgroup 
      \vfill 
      \begingroup
      	\color{Maroon}\spacedallcaps{\myItalianTitle} \\ $\ $\\
      	\spacedallcaps{\myEnglishTitle} \\ 	
	\bigskip
      \endgroup
      \spacedlowsmallcaps{\myName}
      \vfill 
      \vfill
      Relatore: \emph{Relatore}\\
      Correlatore: \emph{Correlatore}\\
      \vfill
      \vfill
      \myTime
      \vfill                      
	\end{center}        
\end{titlepage}   
%--------------------------------------------------------------
% back titlepage
%--------------------------------------------------------------
   \newpage
	\thispagestyle{empty}
	\hfill
	\vfill
	\noindent\myName: 
	\textit{\myItalianTitle,} 
	\myDegree, \textcopyright\ \myTime
%--------------------------------------------------------------
% back titlepage end
%--------------------------------------------------------------
\pagestyle{scrheadings}
%--------------------------------------------------------------
% Mainmatter
%--------------------------------------------------------------
\pagenumbering{arabic}
% use \cleardoublepage here to avoid problems with pdfbookmark
%\chapter{Introduzione}
 % use \myChapter command instead of \chapter
%\tableofcontents\listoffigures\cleardoublepage
\thispagestyle{empty}
\begin{flushright}
\null\vspace{\stretch {1}}
\emph{"Inserire citazione" \break --- Inserire autore citazione} \vspace{\stretch{2}}\null
\end{flushright}
\cleardoublepage
\chapter{Introduzione}

\chapter{Algoritmi di odrinamento e classi di pattern}
In questo capitolo verranno introdotti i principali algoritmi di ordinamento che utilizzano sorting devices, in particolare stack-sort,
queue-sort e bubble sort, e altri concetti necessari per l'analisi della loro composizione.\\\\
Durante l'esecuzione questi algoritmi possono 
salvare gli elementi in un contenitore (la diversa struttura dati
adottata definise i diversi algoritmi) dalla quale poi vengono prelevati
per essere aggiunti all'output.\\\\
Una sola iterazione non garantisce l'ordinamento della permutazione,
dunque gli algoritmi devono essere iterati pi\'u volte, ogni volta sul risultato della iterazione precedente. In ogni caso alla fine delle i-esima
iterazione i maggiori $i$ elementi avranno raggiunto la loro posizione finale, 
dunque sono necessari al massimo $n-1$ iterazioni per ordinare la permutazione.\\\\
Essendo interessati al comportamento di una sola iterazione di questi
algoritmi si esaminer\'a un operatore, definito appositamente per ogni algoritmo, che descrive la singola iterazione. \\
Ad esempio, prendendo l'algoritmo bubble-sort si far\'a riferimento all'operatore $B(\pi)$, dove $\pi$ \'e una permutazione di
interi, tale che n iterazioni del bubble-sort possano essere rappresentate
da $B^n(\pi) = B(\dots B(\pi)\dots)$.
\section*{Bubble sort}
L'algoritmo di ordinamento bubble-sort prevede di scorrere gli elementi da 
ordinare dal primo al penultimo, ed ogni volta confrontare ogni elemento 
con il suo successivo per scambiarli se non sono ordinati.\\
Il risultato di una singola iterazione di bubble-sort su una permutazione $\pi=\pi_1\pi_2\dots\pi_n$ \'e calcolato dall'operatore $B(\pi)$. Per una permutazione $\pi$ con valore massimo $n$ vale che $\pi = \pi_Ln\pi_R$, allora $B(\pi) = B(\pi_L)\pi_rn$.
\begin{algorithm}[H]
   \caption{$B(\pi)$}
\begin{algorithmic}[1]
   \For{$i=1$ {\bfseries to} $n-1$}
   \If{$\pi_i > \pi_{i+1}$} 
   \State Swap $\pi_i$ and $\pi_{i+1}$
   \EndIf
   \EndFor
\end{algorithmic}
\end{algorithm}
\section*{Stack sort}
L'operatore $S(\pi)$ rappresenta il risultato ottenuto applicando un'iterazione di stack sort su una permutazione $\pi$.\\
Il primo passo consiste nell'inserire $\pi_1$ nella pila. Poi lo si confronta con l'elemento $\pi_2$. Se $\pi_1>\pi_2$ allora il secondo viene messo nella pila sopra $\pi_1$, altrimenti $\pi_1$ viene estratto dalla pila e inserito nell'output e $\pi_2$ viene inserito nella pila.\\
Gli stessi passi vengono eseguiti per tutti gli altri elementi presenti nell'input, se viene trovato un elemento nell'input maggiore dell'elemento in cima alla pila, la pila viene svuotata finch\'e questa condizione non diviene falsa, poi l'elemento viene spinto nella pila.\\
Finiti gli elementi nell'input, se necessario, si svuota completamente la pila nell'output.\cite{limbrief}\\
Sia $\pi = \pi_Ln\pi_R$, con $n$ valore massimo in $\pi$, vale che $S(\pi)=S(\pi_L)S(\pi_R)n$
\begin{algorithm}[H]
   \caption{operatore S - stack sort, singola iterazione}
\begin{algorithmic}[1]
\State initialize an empty stack
   \For{$i=1$ {\bfseries to} $n-1$}
   \While{ stack is unempty \textbf{and} $\pi_i>$top of the stack}
   \State pop from the stack to the output
   \EndWhile
   \State push($\pi_i$)
   \EndFor
   \State emtpy the stack in the output
\end{algorithmic}
\end{algorithm}
\section*{Queue sort}
Per ogni elemento $\pi_i$ della permutazione $\pi$ in input se la coda \'e vuota o il suo ultimo elemento \'e minore di $\pi_i$, si accoda $\pi_i$, altrimenti si tolgono elementi dalla coda ponendoli nell'output fino a che l'elemento davanti  alla coda non \'e maggiore di $\pi_i$, poi si aggiunge $\pi_i$ all'output. Si svuota la coda nell'output\cite{magnusson2013sorting}. 
\begin{algorithm}[H]
   \caption{operatore Q - queue sort, singola iterazione}
\begin{algorithmic}[1]
\State initialize an empty queue
\For{$i=1$ {\bfseries to} $n-1$}
\If{empty queue \textbf{or} last in queue $<\pi_i$}
\State enqueue($\pi_i$)
\Else
\While{first in queue $<\pi_i$}
\State dequeue($\pi_i$)
\EndWhile
\State add $\pi_i$ to the output
\EndIf
\EndFor
\State emtpy the queue in the output
\end{algorithmic}
\end{algorithm}
\paragraph*{Bypass} L'operazione che pone un elemento nell'output senza passare dal contenitore si dice \textbf{bypass}. Questa viene svolta normalmente nel queuesort, ma talvolta pu\'o essere introdotta in altri algoritmi.
\paragraph*{Osservazione}\textit{Bubble sort} \'e un caso particolare sia di \textit{queue sort} che di \textit{stack sort}.\\
Se infatti si fissa a 1 la dimensione della pila o della coda dei rispettivi operatori il comportamento che questi assumono \'e quello di una cella che, scorrendo l'input, contiene sempre il massimo valore trovato, mentre gli altri vengono messi nell'output.
\paragraph*{Contenitori POP}Un caso di studio interessante \'e quello in cui i contenitori di stack sort e queue sort vengano sostituiti dalla loro versione POP, cio\'e che quando viene eseguita un estrazione il contenitore viene svuotato completamente. Si definiscono con questa variante, gli algoritmi \textbf{pop-stacksort} e \textbf{pop-queuesort} 
\section*{Classi di pattern di permtuazioni}
\cite{bouvel2022preimages}
Siano $\alpha,\beta$ due sequenze di interi, si indica con $\alpha \subseteq \beta$ che $\alpha$ \'e una sottosequenza di $\beta$, anche se non necessariamente una sottosequenza consecutiva.\\\\
Si dice che una permutazione $\delta$ \'e un pattern contenuto in una permutazione $\tau$ se esiste una sottosequenza di $\tau$ di ordine isomorfico rispetto a $\delta$, e si indica con $\delta\preceq\tau$. Ad esempio $24153$ contiene il pattern $312$ perch\'e $413\subset{24153}$.
\paragraph*{Pattern Standard}
Una \textbf{standardizzazione}\cite{claesson2012sorting} di una sequenza di numeri \'e un altra sequenza della stessa lunghezza in cui l'elemento minore della sequenza originale \'e stato sostituito da 1, il secondo minore con un 2, ecc.\\
Ad esempio, la standardizzazione di 5371 \'e 3241.
\paragraph*{Classi di pattern} La relazione di sottopermutazione \'e una relazione di ordine parziale che viene studiata con dei sottoinsiemi chiamati \textbf{pattern di classi}. Ogni classe di pattern $D$ pu\'o essere caratterizzata dall'insieme minimo $M$ che evita:$$ D = Av(M) = \{\beta:\mu\not\preceq\beta\forall\mu\in M\}$$
\section*{Algoritmi e pattern-avoidance}
\paragraph*{pattern 231}Una permutazione $\pi$ contiene un pattern 231 se $231\preceq\pi$, ovvero se $\exists a<b<c: \pi =\dots{b}\dots{c}\dots{a}\dots$.\\\\
\'E noto in letteratura \cite{limbrief} che una permutazione pu\'o essere ordinata da una sola passata di stack sort se e solo se non contiene pattern 231.\\\\
Allo stesso modo, sono note simile condizioni perch\'e una permutazione possa essere ordinata da una sola passata degli altri operatori descritti precedentemente.
\begin{center}
\begin{tabular}{ |c|c| } 
\hline
\textbf{Operatore} & \textbf{Permutazioni ordinabili con una sola passata} \\ 
\hline
Stack sort & Av(231)\\ 
Queue sort & Av(321)\\ 
Bubble sort & Av(231,321)\\ 
Pop-stack sort & Av(231,312)\\ 
Pop-queue sort& Av(132,321)\\ 
\hline
\end{tabular}
\end{center}

\chapter{Programmi realizzati}
In questo capitolo presenter\'o il funzionamento di alcuni programmi che ho realizzato come supporto allo studio degli argomenti di questa tesi.\\\\
Tutto il software che verr\'a presentato \'e stato realizzato con il linguaggio Python, nel seguente ambiente: \texttt{3.11.6 (main) [GCC 13.2.1 20230801]}. Tutti gli script vengono lanciati da linea di comando e mostrano i risultati sulla console.
\section*{PermutaSort}
Il primo programma che ho realizzato viene lanciato da linea di comando specificando come argomenti un intero $n$ e un operatore di ordinamento $X$. Il suo scopo \'e quello di enumerare tutte le $n$-permutazioni ordinabili e non ordinabili con una sola passata di $X$, stampando anche i possibili risultati di $X$ su $n$-permutazioni.\\\\
La classe \texttt{SelectorPermutations} \'e il core del programma, viene inizializzata passando un intero e un operatore al costruttore. Grazie alla funzione \texttt{itertools.permutations} vengono generate le $n$-permutazioni e , scorrendole tutte, vi si applica l'operatore $X$. Il risultato viene aggiunto alla lista \texttt{outcomes} e viene valutato: se \'e ordinato la permutazione da cui si \'e ottenuto viene aggiunta alla lista \texttt{sortable}, altrimenti a \texttt{unsortable}.\\La classe fornisce dei \textit{getters} per le liste e non ha altri metodi, oltre al costruttore.\\
Dopo essere stata istanziata la classe viene usata come riferimento per i dati che ha gi\'a calcolato nel costruttore, e non produce ulteriori risultati.\\Di seguito viene mostrato il listato della classe:
\\\lstinputlisting[language=Python]{./python/permutasort.py}
%\pagebreak
\section*{Pattfinder}
Questo programma ha lo scopo di selezionare tra tutte le permutazioni di una data dimensione quali contengono un dato pattern classico e quali no.\\\\
Il programma viene lanciato da linea di comando prendendo come argomenti un intero e una sequenza di interi consecutivi eventualmente non ordinati.\\
Il primo intero rappresenta la lunghezza delle permutazioni da scorrere e la sequenza rappresenta il pattern da ricercare.\\
Gli argomenti vengono passati al costruttore della classe \texttt{PatternAvoid}, l'intero $n$ viene utilizzato per generare le $n$-permutazioni grazie alla funzione \texttt{permutations}. Per ogni permutazione viene controllato se contiene o no il pattern e viene inserita nella rispettiva lista: \texttt{containing} o \texttt{notcontaining}.\\\\
La classe presenta i getter per le liste prodotte e due ulteriori metodi: 
\begin{description}
\item[\texttt{patternize}] serve a \textit{standardizzare} una sequenza, si fa riferimento alla definizione di standardizzazione (di una permutazione) presentata nel capitolo precedente
\item[\texttt{contains}] verifica che una sequenza contenga un pattern o meno, grazie alla funzione \texttt{itertools.combinations} si ottengono tutte le sottosequenze (anche non consecutive), le si standardizzano e le si confrontano con il pattern da ricercare.
\end{description}
\lstinputlisting[language=Python]{./python/pattfinder.py}
\section*{Script per la verifica dei pattern}
Spesso si \'e utilizzato un ulteriore script, che si serve delle classi presentate prima, per verificare che i pattern trovati per un dato operatore siano sufficienti a definire tutto l'insieme di permutazioni non ordinabili.\\\\
Grazie a \texttt{selectorPermutations} vengono generate tutte le permutazioni non ordinabili (l'operatore e la dimensione delle permutazioni vengono passate come argomenti allo script) ed in seguito si scorre la lista di pattern che si vogliono veririficare. Per ogni pattern \texttt{PatternAvoid} genera le permutazioni che lo contengono e queste vengono rimosse dalla lista di permutazioni non ordinabili. Se alla fine la lista rimane vuota allora viene stampato a schermo che tutte le permutazioni non ordinabili contengono almeno uno dei pattern specificati, altrimenti elenca quelle rimaste.\\\\
Lo script viene lanciato da linea di comando con i seguenti argomenti (nell'ordine): l'operatore di ordinamento, la dimensione delle permutazioni, l'elenco di pattern da verificare.\\\\
Di seguito si presenta un esempio di utilizzo ed il codice dello script; l'utilizzo mostrato \'e in riferimento alla combinazione degli operatori stacksort e bubblesort, che verr\'a speigato mostrata nel dettaglio nel capitolo successivo.\\
\lstinputlisting[language=bash]{./python/result_verify.txt}
\lstinputlisting[language=Python]{./python/test.py}

\chapter{Composizione di operatori di ordinamento}
\'E molto interessante studiare la combinazione dei vari operatori e le relazioni tra classi di pattern e le loro controimmagini.
\paragraph*{controimmagini} Sia $X$ un operatore di ordinamento, la controimmagine di una certa permutazione $p$, secondo $X$, indicata con $X^{-1}(p)$ rappresenta l'insieme di tutte le permutazioni la cui immagine secondo l'operatore $X$ \'e uguale a $p$.$$X^{-1}(p) = \{\beta : p = X(\beta)\}$$\\
$Av(21)$ \'e l'insieme formato dalle sole permutazioni identit\'a, dato che contiene tutte le permutazioni in cui nessun elemento sia disordinato rispetto ad un altro, ovvero solo le permutazioni crescenti. Si indica con $X^{-1}(Av(21))$ l'insieme di tutte le permutazioni ordinabili da $X$.\\\\
Ad esempio, dato che \'e noto che l'operatore \textit{bubblesort} ordina solo le permutazioni che non contengono pattern 231 e 321\cite{claesson2012sorting}, vale che:$$B^{-1}(Av(21)) = Av(231,321)$$
\paragraph*{Composizione di operatori} Siano due operatori di ordinamento $X$ e $Y$, la loro composizione \'e indicata con $ XY = X \circ Y$, qundi, ad esempio, la composizione di \textit{Stacksort} e \textit{Bubblesort} (in questo ordine) si indica con $SB(\pi) = (S \circ B)(\pi) = S(B(\pi))$.
\paragraph*{Algoritmi per il calcolo di controimmagini} Sono gi\'a stati prodotti alcuni algoritmi per calcolare le controimmagini di pattern secondo gli operaratori \textit{bubblesort}\cite{albert2010inverse}, \textit{stacksort}\cite{claesson2012sorting} e \textit{queuesort}\cite{magnusson2013sorting}\cite{cioni2021characterization}.\\
Questo ci permette, quando si combinano due operatori di ordinamento (di cui di quello applicato per secondo siano noti le classi di pattern da evitare per l'ordinabilit\'a) di cercare per quali pattern l'operatore che viene applicato per primo produce permutazioni che siano ordinabili dal secondo.
\paragraph*{Un caso gi\'a noto: Combinazione di \textit{Stacksort} con \textit{Bubblesort}} Si considera dunque la composione $SB = S\circ{B}$ e ci si chiede quali permutazioni possano essere ordinate da esso.$$(SB)^{-1}Av(21)=B^{-1}S^{-1}(Av(21))=B^{-1}(Av(231))$$
Dunque ricercando per quali permutazioni \textit{bubblesort} evita il pattern 231 si trovano le condizioni per cui una permutazione risulta ordinabile da \textit{SB}.\\
Utilizzando l'algoritmo per le controimmagini di bubblesort\cite{albert2010inverse} si ottiene che:$$(SB)^{-1}(Av(21))=Av(3241, 2341, 4231, 2431)$$
Per ottenere una controimmagine di bubblesort di un pattern classico, si applicano le seguenti regole:
\begin{description}
	\item si cercano tutte le inversioni nell'immagine in esame: poich\'e bubblesort non produce nuove inversioni, tutte le inversioni nell'immagine dovranno gi\'a essere presenti nella controimmagine;
	\item si considera una lista di \textbf{candidati}, composta da ogni pattern minimale che contiene almeno le stesse inversioni;
	\item per ogni candidato, si processa ogni sua coppia disordinata $(b,a)$ nel modo seguente:
	\begin{description}
	\item se $(b,a)$ \'e contenuta anche nell'immagine si ha che $b$ non \'e un LTR massimo o che lo \'e ma c'\'e un altro LTR massimo tra $b$ e $a$, quindi si decora la coppia come segue, producendo un nuovo pattern:
	\begin{center}
		\mmpattern{scale=1.4}{2}{1/2,2/1}{}{}$\Rightarrow$\mmpattern{scale=1.4}{2}{1/2,2/1}{}{0/2/2/3/1}
	\end{center}
	\item se $(b,a)$ \'e un'inversione del candidato che non \'e nell'immagine allora la coppia viene ordinata da una passata di bubblesort, si ha che $b$ \'e l'ultimo LTR massimo prima di $a$, la coppia viene decorata nel modo seguente:
	\begin{center}
		\mmpattern{scale=1.4}{2}{1/2,2/1}{}{}$\Rightarrow$\mmpattern{scale=1.4}{2}{1/2,2/1}{0/2,1/2}{}\\
	\end{center}
	\end{description}
	\item applicare ad un candidato le condizioni per tutte le sue inversioni (considerando che se una zona \'e oscurata da una condizione, non pu\'o essere decorata da un altra) fornisce un mesh-pattern decorato che descrive il pattern contenuto nella controimmaigne.
\end{description}
\subsection*{Esempio pratico di calcolo di una controimmagine}
A titolo d'esempio verr\'a mostrato come si pu\'o arrivare al risultato (gi\'a noto\cite{albert2010inverse}) riguardo alla combinazione $(S\circ B)$:
$$(SB)^{-1}(Av(21)) = Av(2341, 2431, 3241, 4231)$$
Prima di tutto si ricerca di quale classe di pattern nonch\'e tramite quale operatore occorre calcolare le controimmagini:
$$(SB)^{-1}(Av(21)) = B^{-1}(S^{-1}(Av(21)))=B^{-1}(Av(231))$$
Dato che, come noto\cite{limbrief} $S^{-1}(Av(21))=Av(231)$.\\
I pattern candidati ad essere controimmagini di 231 sono i pattern che contengono le coppie $(2,1),(3,1)$, ovvero i pattern 231,321.
\paragraph*{231:} si osserva che una condizione per l'inversione $(2,1)$ \'e gi\'a realizzata dall'elemento $3$, applicando le regole descritte prima all'inversione $(3,1)$ si ha il mesh-pattern \mmpattern{scale=0.7}{3}{1/2,2/3,3/1}{}{0/3/3/4/1}, ovvero i pattern 2341, 2431, 4231.
\paragraph*{321:} in questo caso si ha l'inversione $(3,2)$ che deve essere invertita, quindi si inseriscono le aree oscurate, in seguito si osserva che, come prima, $3$ realizza la condzione per l'inversione $(2,1)$ e dunque basta aggiungere le decorazioni per l'inversione $(3,1)$. Di seguito si mostrano i passaggi:
\begin{center}
\mmpattern{scale=1.5}{3}{1/3,2/2,3/1}{}{}$\Rightarrow$
\mmpattern{scale=1.5}{3}{1/3,2/2,3/1}{0/3,1/3}{}$\Rightarrow$
\mmpattern{scale=1.5}{3}{1/3,2/2,3/1}{0/3,1/3}{2/3/3/4/1}\\
\end{center}
Si ottiene cos\'i il pattern 3241 che, unito agli altri fornisce esattamente il risultato atteso.\\\\
Si osservi che il software \texttt{permutasort}, presentato nel capitolo 3, pu\'o aiutare ad arrivare alle stesse conclusioni. Infatti lanciando il programma sull'operatore SB e con $n=4$ si ottengono i seguenti risultati:\\\\\texttt{\$ python permutasort.py 4 SB }\\\dots\\\texttt{The following 4 4-permutations are not sortable with the operator SB:}\\\texttt{(2, 3, 4, 1)}\\\texttt{(2, 4, 3, 1)}\\\texttt{(3, 2, 4, 1)}\\\texttt{(4, 2, 3, 1)}\\\dots\\\\
Questo tuttavia \'e utile come indicazione, ma non \'e sufficiente per essere sicuri del risultato: infatti non si pu\'o escludere la presenza di pattern classici pi\'u lunghi n\'e la presenza di pattern barrati.\\
Per quanto quindi questo strumento si sia rivelato molto utile per avere un'idea di quali pattern ricercare, soprattutto durante le combinazioni pi\'u complesse che presenteremo in seguito, non pu\'o sostituire del tutto un approccio pi\'u formale e teorico, come l'applicazione di uno degli algoritmi introdotti prima.
\section*{Combinazione di queuesort con bubblesort}
$$(QB)^{-1}(Av(21))=B^{-1}Q^{-1}(Av(21))=B^{-1}(Av(321))$$
Si osserva che le inversioni in $321$ sono $(3,2),(3,1),(2,1)$ e l'unico pattern minimo che le contiene tutte \'e appunto $321$, che quindi \'e l'unico candidato.\\
In questo caso a tutte le inversioni deve essere applicata la prima condizione, a parte a $(2,1)$ perch\'e la condzione \'e gi\'a soddisfatta da 3. Il pattern decorsato che si ottiene \'e:
\begin{center}
\mmpattern{scale=1.5}{3}{1/3,2/2,3/1}{}{0/3/3/4/{\dots\dots1},0/3/2/4/1}\\
\end{center}
Ovvero  l'unione di \mmpattern{scale=0.7}{3}{1/3,2/2,3/1}{}{0/3/2/4/1} e \mmpattern{scale=0.7}{3}{1/3,2/2,3/1}{}{0/3/3/4/1}. Tuttavia questi due pattern possono essere semplificati nel primo: si nota infatti che \mmpattern{scale=0.7}{3}{1/3,2/2,3/1}{}{0/3/2/4/1}, che si ottiene dalle condizioni per $(3,2)$ soddisfa anche le condizioni per $(3,1)$ e inoltre vale che \mmpattern{scale=0.7}{3}{1/3,2/2,3/1}{}{0/3/2/4/1}$\preceq$\mmpattern{scale=0.7}{3}{1/3,2/2,3/1}{}{0/3/2/4/1,2/3/3/4/1}.\\
Quindi le preimmagini del pattern 321 tramite bubblesort sono date da:
\begin{center}
\mmpattern{scale=1.4}{3}{1/3,2/2,3/1}{}{0/3/2/4/1}.
\end{center}
Questo mesh pattern corrisponde ai pattern classici $4321,3421$.$$QB^{-1}(Av(21))=Av(4321,3421)$$
\section*{Composizioni di un operatore di ordinamento con Stacksort}
L'algoritmo per le controimmagini di \textit{stacksort} \'e abbastanza simile a quello per \textit{bubblesort}.\\
Anche in questo caso si osservano tutte le inversioni contenute nell'immagine in esame e se ne trae una lista di pattern minimali candidati ad essere controimmagini e li si esaminano ad uno ad uno.\\
Per ogni coppia disordinata $(b,a)$ presente nel candidato: 
\begin{description}
\item se $(b,a)$ \'e presente anche nell'immagine allora deve essere presente un elemento $c>b$ tra $b$ e $a$ che fa uscire $b$ dalla pila prima che $a$ vi entri
\begin{center}
\mmpattern{scale=1}{2}{1/2,2/1}{}{} $\Longrightarrow$\mmpattern{scale=1}{2}{1/2,2/1}{}{1/2/2/3/1} 
\end{center}
\item se invece $(b,a)$ non \'e presente allora si pu\'o escludere la presenza di tale elemento
\begin{center}
\mmpattern{scale=1}{2}{1/2,2/1}{}{} $\Longrightarrow$\mmpattern{scale=1}{2}{1/2,2/1}{1/2}{} 
\end{center}
\end{description}
\subsection*{Composizione  di {queuesort} con {stacksort}}
$$(QS)^{-1}(Av(21))=S^{-1}Q^{-1}(Av(21))=S^{-1}(Av(321))$$
Si ricercano dunque le controimmagini di $321$ secondo \textit{stacksort}.\\$321$ \'e l'unico candidato, quindi sappiamo che tutte le controimmagini devono contenere il pattern $321$.\\
L'elemento $3$ deve entrare nella pila ed uscirne prima del $2$, per fare ci\'o deve esserci un elemento $3^+$ maggiore di $3$ tra $3$ e $2$. Similmente \'e necessario un elemento $2^+>2$ tra $2$ e $1$ per assicurare che $2$ esca dalla pila prima che $1$ vi entri.\\
\begin{center}
\mmpattern{scale=1.5}{3}{1/3,2/2,3/1}{}{1/3/2/4/1,2/2/3/4/1}
\end{center}
Dunque le controimmagini che cerchiamo devono essere nella forma $33^+22^+1$.\\
Se $2^+<3$ allora la controimmagine assume la forma del pattern $45231$.\\
Altrimenti se $2^+>3$, si possono ottenere due diverse controimmagini:
\begin{description}
	\item[$3^+>2^+$] genera il pattern $35241$
	\item[$2^+>3^+$] genera il pattern $34251$
\end{description}
$$(QS)^{-1}(Av(21)) = Av(34251, 35241, 45231)$$
\subsection*{Composizione di {bubblesort} con {stacksort}}
$$(BS)^{-1}(Av(21))=S^{-1}B^{-1}(Av(21))=S^{-1}(Av(231,321))$$
Gi\'a dall'analisi della combinazione precedente \'e risultato che $S^{-1}(Av(321))=Av(34251, 35241, 45231)$ qundi \'e necessario calcolare solo $S^{-1}(Av(231))$. Quest'ultimo risultato \'e stato ampiamente studiato in letteratura, in quanto analogo al caso di una variante di \textit{stacksort} che utilizza 2 pile. Il risultato che si ottiene \'e dunque che $S^{-1}(Av(231))=Av(2341, 3\overline{5}241)$\cite{claesson2012sorting}.\\\\
Seguendo l'algoritmo applicato finora si osserva che i pattern candidati per le controimmagini sono $231,321$.\\\\
Si uniscono dunque i due risultati:
$$S^{-1}(Av(231))=Av(2341,3\overline{5}241), S^{-1}(Av(321))=Av(34251, 35241, 45231)$$
Si osserva che $2341\preceq 34251$, quindi $34251$ non \'e minimo.\\
Inoltre i pattern $35241, 3\overline{5}241$, possono essere rappresentati dal pattern minimo $3241$ che rende non minimo anche $34251$.\\Il risultato che si ottiene \'e:$$(BS)^{-1}(Av(21))=Av(2341,3241,45231)$$
\section*{Combinazioni che terminano con \textit{queuesort}}
Nonostante anche per \textit{queuesort} sia stato trovato un algoritmo per le controimmagini\cite{magnusson2013sorting}, approcciare il problema "manualmente" analizzando i possibili comportamenti di \textit{queuesort} rispetto alle diverse possibili permutazioni risulta essere pi\'u semplice e pi\'u comprensibile.\\\\
Se nel valutare il comportamento di \textit{queuesort} rispetto ad una controimmagine se ci si vuole assicurare che un elemento entri nella coda allora lo si vuole "forzare" ad essere un LTR maxima, dunque si applicher\'a il seguente mesh-pattern all'elemento in considerazione:
\begin{center}
\mmpattern{scale=1.5}{1}{1/1}{0/1}{}
\end{center}
Al contrario, se si ha a che fare con un elemento che risulta essere un right-to-lrft-maxima ma lo si vuole forzare a svuotare la coda, si avr\'a il seguente risultato:
\begin{center}
\mmpattern{scale=1.5}{1}{1/1}{}{0/1/1/2/1}
\end{center}
Per forzare un elemento a effettuare un bypass \'e necessario che esso sia minore del primo elemento della coda, per questa condizione \'e pi\'u difficile formalizzare una startegia, ma si vedranno in seguito esempi di come pu\'o essere fatto. 
\subsection*{Composizione di {stacksort} con {queuesort}}
$$(SQ)^{-1}(Av(21))=Q^{-1}S^{-1}(Av(21))=Q^{-1}(Av(231))$$Ancora una volta si considerano come candidati i pattern minimali che contengano almeno le inversioni $(2,1),(3,1)$ ovvero $231,321$.\\\\
Se gli elementi del pattern $231$ nell'immagine formano un pattern $321$ nella controimmagine allora il valore $3$ deve entrare nella coda. Sia che $2$ effettui un bypass sia che provochi l'estrazione di alcuni elementi della coda comunque esso viene aggiunto all'output prima che $3$ venga estratto. Si rende necessario un elemento $3^+>3$ posto tra $2$ e $1$ per estarre il $3$ dalla coda e inserirlo nell'input prima dell'$1$, tuttavia $3^+$ non deve essere maggiore dell'elemento in fondo alla coda, altrimenti viene accodato. Si teorizza l'esistenza di un elemento $4>3^+$ posto tra $3$ e $2$, che rappresenta il massimo della coda quando si raggiunge $3^+$.\\
Il pattern che si ottiene deve essere nella forma $3423^+1$, ovvero $35241$\\\\
se invece gli elementi del pattern $231$ nell'immagine formano un pattern $231$ anche nella controimmagine che lo contiene a sua volta \'e necessario che $2$ effettui un bypass o che venga estratto dalla coda prima che $3$ vi entri. Il che implica la presenza di un valore $2^+>2$ prima di $2$ o tra $2$ e $3$.\\Similmente $3$ deve essere aggiunto all'output prima di $1$, qundi deve bypassare o essere estratto prima che $1$ bypassi. Il che implica l'esistenza di un elemento $4$ in posizione precdente a $1$; in particolare i pattern minimali che rispettano queste condizoni sono quelli in cui $4$ compare prima di $2$ e si ha che $2^+=4$.\\
\begin{center}
\mmpattern{scale=1.5}{3}{1/2,2/3,3/1}{}{0/3/2/4/1}
\end{center}
Gli unici pattern minimali che soddisfano queste condizioni sono $2431,4231$.
\\\\Si nota inoltre che, il pattern calcolato prima $35241\in Q^{-1}(Av(321))$ risulta adesso superfluo in quanto $4231\preceq35241$
$$(SQ)^{-1}(Av(21))=Av(2431,4231)$$
\subsection*{Composizione di {bubblesort} con {queuesort}}
$$(BQ)^{-1}(Av(21))=Q^{-1}B^{-1}(Av(21))=Q^{-1}(Av(231,321))$$
Ancora una volta, si osserva che $Q^{-1}(Av(231))$ \'e gi\'a stato calcolato nel caso precedente, e resta dunque da unire i risultati gi\'a ottenuti con $Q^{-1}(Av(321))$.\\
Gli elementi del pattern $321$ nell'immagine devono essere a sua volta nello stesso odrine nell'immagine.\\
In particolare gli elementi $3,2$ devono effettuare un bypass o essere inseriti nella coda ed estratti prima dell'elemento successivo.\\\\
\begin{center}
\mmpattern{scale=1.5}{3}{1/3,2/2,3/1}{}{0/3/2/4/1}
\end{center}
Il pattern $4321$ sicuramente soddisfa queste condizioni: se $4$ bypassa allora tutti gli elementi del pattern, essendo minori bypassano a loro volta, se viene inserito in coda ed estratto allora bypassano comunque, se viene inserito in coda e non estratto allora tutti gli elementi successivi evitano la coda e vengono aggiunti all'output, eventualmente dopo aver provocato l'estrazione di alcuni elementi dalla coda.\\
Se $3$ viene aggiunto in coda ed estratto allora si ha una situazione $3421$, dove $2$ bypassa.\\
$$(SQ)^{-1}(Av(21))=Av(2431,3421,4231,4321)$$
\section*{Contenitori POP}
Lo studio di combinazioni di operatori che utilizzano contenitori POP si rivela pi\'u difficile, in quando si tratta di casi non approfonditi in letteratura.\\
Sono tuttavia noti i pattern che rendono una permutazione non ordinabile, come presentato nel capitolo 2, ovvero:
$$S_{POP}^{-1}(Av(21)) = Av(231,312)$$$$Q_{POP}^{-1}(Av(21))=Av(321,2413)$$
Grazie a questi dati, \'e possibile individuare un caso di facile risoluzione: ovvero quello in cui un algoritmo che usa contenitori POP \'e concatenato ad un operatore regolare.
$$(X_{POP}\circ{Y})(\pi)= X_{POP}(Y(\pi)))\Rightarrow(X_{POP}\circ{Y})^{-1}(\pi) = Y^{-1}X_{POP}^{-1}(\pi)=Y^{-1}(Av(m))$$
Sapendo l'insieme di pattern $m$ che rendono la permutazione non ordinabile dall'operatore POP baster\'a cercare le loro controimmagini secondo l'operatore regolare con i metodi usati finora, come verr\'a mostrato nelle sezioni seguenti.
\section*{POP Queuesort}L'algoritmo POP queuesort richiede una diversa analisi da POP stacksort.\\
Esistono diverse versioni di \textit{POP queuesort} ed in particolare ne esistono due ottimali\cite{cioni2021sorting}: \textit{Min} e \textit{Cons}.
\begin{description}
	\item[\textit{Min}] in questa versione l'operazione di POP viene eseguita solo se il primo elemento della coda \'e il successivo dell'ultimo elemento aggiunto all'output; se l'elemento in input \'e maggiore dell'ultimo elemento della coda (o se la coda \'e vuota) allora viene accodato mentre negli altri casi, se l'elemento dell'input \'e minore della testa della coda allora bypassa altrimenti si esegue un POP.
	\item[\textit{Cons}] questa versione si basa sull'idea di avere sempre elementi consecutivi nella coda; \'e la versione che verr\'a adottata in questa tesi; da qui in avanti ogni riferimento a \textit{POP-queuesort} sar\'a riferito a \textit{Cons}
\end{description}
\begin{algorithm}[H]
   \caption{Cons - POP Queuesort}
\begin{algorithmic}[1]
\State $Q\leftarrow$ empty POP-queue
\For{$i=1$ {\bfseries to} $n-1$}
	\If{empty queue \textbf{or} $\pi_i=Back(Q)+1$}
		\State enqueue($\pi_i$)
	\Else
		\If{ $Front(Q)>\pi_i$}
			\State append $\pi_i$ to the output
		\Else
			\State POP the queue and append the result to the output
			\State enqueue($\pi_i$)
		\EndIf
	\EndIf
\EndFor
\State POP the queue and append the result to the output
\end{algorithmic}
\end{algorithm}
La dimostrazione che \textit{Cons} (cos\'i come \textit{Min}) sia un algoritmo ottimale nella classe degli algoritmi \textit{POP-queueosrt} si ha dal fatto che esso ordina tutte e sole le sequenze dell'insieme Av(321, 2413), che sono esattamente tutte e sole le permutazioni ordinabili con una POP queue\cite{cioni2021sorting}.
\subsection*{Composizione di {POP-stacksort} con {stacksort}}
$$(S_{POP}\circ{S})^{-1}(Av(21))=S^{-1}(Av(231,312))$$
Si \'e gi\'a calcolato che $S^{-1}(Av(231))=Av(2341, 3\overline{5}241)$, quindi manca da calcolare le controimmagini di $312$.\\
I pattern candidati che possono generare $312$ sono quelli che contengono le coppie $(3,1),(3,2)$, ovvero $321, 312$.
Perch\'e gli elementi di un pattern $321$ nella controimmagine generino un pattern $312$ tramite \textit{stacksort} \'e necessario che $3$ entri in pila e ne esca prima che $2$ vi entri, deve quindi essere presente un valore $4$ posizionato tra $3$ e $2$ che provochi l'uscita di $3$, ovvero il pattern $3421$.\\
Se gli elementi del pattern $312$ nell'immagine formano lo stesso pattern anche nella controimmagine allora anche in questo caso bisogna assicurarsi che $3$ venga estratto dalla pila prima di $1$, si ha quindi il pattern $3412$.\\Si uniscono tutti i risultati:
$$(S_{POP}S)^{-1}(Av(21))=Av(2341, 3412, 3421, 3\overline{5}241)$$
\subsection*{Composizione di {POP-stacksort} con {bubblesort}}
$$(S_{POP}\circ{B})^{-1}(Av(21))=B^{-1}(Av(231,312))$$
Anche in questo caso la controimmagine di $231$ \'e stata calcolata: $B^{-1}(Av(231)) = (S\circ{B})^{-1}(Av(21)) = Av(2341, 2431, 3241, 4231)$.\\
Si ricerca dunque le preimmaigini di 312 secondo bubblesort. \\
Gli elementi del pattern $312$ nell'immagine possono solo formare un pattern $312$ o $321$ nella controimmagine.\\
Si pu\'o osservare come gli stessi elementi di un pattern $321$ non possono divenire un $312$. La presenza dell'elemento maggiore del pattern all'inizio di essa evita l'ordinamento degli elementi successivi.\\
Perch\'e un pattern $312$ rimanga invariato \'e sufficiente la presenza di un elemento maggiore di tutto il pattern in posizione tale da evitare l'ordinamento di $3$ con qualsiasi elemento, applicando le regole descritte prima per \textit{bubblesort}: 
\begin{center}
\mmpattern{scale=1.5}{3}{1/3,2/1,3/2}{}{0/3/2/4/1}$=4312,3412$
\end{center}
$$(S_{POP}{B})^{-1}(Av(21))=Av(2341, 2431, 3241, 3412, 4231, 4312)$$
\subsection*{Composizione di {POP-stacksort} con {queuesort}}
$$(S_{POP}\circ{Q})^{-1}(Av(21))=Q^{-1}(Av(231,312))$$$$Q^{-1}(Av(231)=Av(4231,2431)$$
Possibili controimmagini di $312: 312, 321$.\\
Se la controimmagine contiene $312$ bisogna far si che $3$ effettui un bypass o che esca prima di $1$, in ogni caso queste condizioni sono soddisfatte dal seguente mesh-pattern, $4312$:
\begin{center}
\mmpattern{scale=1.5}{3}{1/3,2/1,3/2}{}{1/2/2/3/1}
\end{center}
$$(S_{POP}{Q})^{-1}(Av(21))=Av(2431,4231,4312)$$
\subsection*{Composizione di {POP-queuesort} con {stacksort}}
$$(Q_{POP}{S})^{-1}(Av(21)) = S^{-1}(Q_{POP}^{-1}(Av(21))) = S^{-1}(Av(321, 2413))$$
Come in molti altri casi la controimmagine di 321 \'e stata calcolata:$$S^{-1}(Av(321)) = (QS)^{-1}(Av(21)) = Av(34251, 35241, 45231)$$
Si ricerca dunque le controimmagini di 2413 secondo stacksort, applicando lo stesso algoritmo dei casi precedenti.\\
Le inversioni di 2413 sono $(2,1),(4,1),(4,3)$, che definiscono l'insieme di candidati 2413, 2431, 4213, 4231.
\paragraph*{2431} I primi due elementi del pattern sono ordinati fra loro, dunque per lasciare il pattern invariato \'e sufficiente evitare che le coppie 41 e 43 vengano ordinate. Come gi\'a visto questo si ottiene posizionando in mezzo ai due elementi della coppia un elemnto maggiore di entrambi. Si ottiene cos\'i  il pattern 24513.
\paragraph*{2431} Le inversioni di questa controimmagine sono quelle gi\'a individute nell'immagine pi\'u $(3,1)$. Quindi applicando le condizioni per le controimmagini si ottiene il seguente meta pattern:
\begin{center}\mmpattern{scale=1}{4}{1/2,2/4,3/3,4/1}{}{}$\Rightarrow$\mmpattern{scale=1}{4}{1/2,2/4,3/3,4/1}{3/3,3/4}{2/4/3/5/1}\end{center}
La controimmagine di 2413 che contiene 2431 \'e quindi 24531.
\paragraph*{4213}\begin{center}\mmpattern{scale=1}{4}{1/4,2/2,3/1,4/3}{}{}$\Rightarrow$\mmpattern{scale=1}{4}{1/4,2/2,3/1,4/3}{1/4}{2/4/3/5/1}\\La controimmagine di 2413 che contiene 4213 \'e 42513.\end{center}
\paragraph*{4231}\begin{center}\mmpattern{scale=1}{4}{1/4,2/2,3/3,4/1}{}{}$\Rightarrow$\mmpattern{scale=1}{4}{1/4,2/2,3/3,4/1}{1/4,3/3,3/4}{2/4/3/5/1}\\La controimmagine di 4213 \'e 42531.\end{center}
Unendo tutti i risultati si ottiene:$$(Q_{POP}\circ{S})^{-1}(Av(21)) = Av(24513, 24531, 34251, 35241, 42513, 42531, 45231)$$
\subsection*{Composizione di {POP-queuesort} con {bubblesort}}
$$(Q_{POP}{B})^{-1}(Av(21)) = B^{-1}Q_{POP}^{-1}(Av(21))=B^{-1}(Av(321,2413))$$
Come in tutti gli altri casi POP, una delle condizioni \'e gi\'a stata calcolata: $B^{-1}(Av(321))=(Q\circ{B})^{-1}(Av(21))=Av(3421,4321)$. Si ricerca adesso le controimmagini di 2413 (secondo bubblesort).\\
I candidati per le controimmagini sono: 2413, 2431, 4213, 4231, 4321.\\
Si osserva subito che 2431, 4231, e 4321, pur essendo nell'elenco dei candidati non possono in realt\'a generare un pattern 2413 con gli stessi elementi: in ognuno di questi infatti la coppia $(3,1)$ dovrebbe essere ordinata ma nessuno dei due \'e un LTR maxima, dunque il loro ordine relativo non viene cambiato da bubblesort.\\
Si applicano le condizioni per la ricerca di controimmagini di bubblesort agli altri due candidati:
\begin{description}
	\item[2413:]\begin{center}\mmpattern{scale=1}{4}{1/2,2/4,3/1,4/3}{}{}$\Rightarrow$\mmpattern{scale=1}{4}{1/2,2/4,3/1,4/3}{}{0/4/3/5/1}\end{center}
	Dunque i pattern che contengono un pattern 2413 che resta invariato dopo una passata di bubblesort sono 24513, 25413, 52413.
	\item[4213]\begin{center}\mmpattern{scale=1}{4}{1/4,2/2,3/1,4/3}{}{}$\Rightarrow$\mmpattern{scale=1}{4}{1/4,2/2,3/1,4/3}{0/4}{2/4/3/5/1}\end{center}
	Il pattern minimo che corrisponde al mesh pattern rappresentato \'e 42513.\\
	Si noti che l'area oscurata dovrebbe generare i pattern $\overline{6}42513$ e $\overline{5}42613$, tuttavia si nota che entrambi i pattern, se inclusa la parte barrat, sono contenuti in quello precedente e generano pattern 2413 dopo una passata di bubblesort: $B(642513) = B(542613) =425136$, dove la sottosequenza 2513 rappresenta un'istanza del pattern 2413.
\end{description}
\subsection*{Composizione di {POP-queuesort} con {queuesort}}
$$(Q_{POP}{Q})^{-1}(Av(21)) = Q^{-1}(321, 2413)$$
La controimmagine $Q^{-1}(Av(321))$ \'e gi\'a stata trovato nel caso $B\circ{Q}$, e vale che $Q^{-1}(Av(321)) = Av(4321)$.\\
Se si applica le regole per la ricerca di preimamgini di queuesort a 2413 si ottiene:
\begin{center}\mmpattern{scale=1.5}{4}{1/2,2/4,3/1,4/3}{}{0/4/2/5/1}\end{center}
Quindi si ottiene che $(Q_{POP}\circ{Q})^{-1}(Av(21)) = Av(4321, 25413, 52413)$.
\section*{Idee per la ricerca di controimmagini secondo l'operatore POPstacksort}
Si ricercano le condizioni che un candidato deve soddisfare per produrre una voluta immagine. Come negli altri casi, si esaminano le inversioni nel candidato, per ogni coppia disordinata $(b,a)$ che non viene ordinata da una passata di POPstacksort si verifica una delle sequenti condizioni:
\begin{description}
	\item[a)] un elemento $c>b$ \'e presente tra $b$ e $a$, questo provoca il POP prima che l'algoritmo arrivi ad esaminare $a$;
	\item[b)] un elemento $a^-<a$ (e nessun elemento $c>b$) \'e presente tra $b$ e $a$, quindi $a^-$ viene inserito in coda ed \'e $a$ a provocare il POP;
	\item[c)] tra i due elementi ma non c'\'e nessun elemento minore di $a$ o maggiore $b$ ma c'\'e una coppia ordinata, ovvero esistono ${b^-,a^+}:a<a^+<b^-<b$ e $b\dots{a^+}\dots{b^-}\dots{a}$ \'e una sottosequenza nel candidato.
\end{description}
Quindi per trovare una controimmagine secondo POPstacksort occorre applicare le seguenti decorazioni, per ogni coppia disordinata che debba restare tale, al meta pattern del candidato:
\begin{center}
\mmpattern{scale=1.5}{2}{1/2,2/1}{}{}$\Rightarrow$
A)\mmpattern{scale=1.5}{2}{1/2,2/1}{}{1/2/2/3/1},
B)\mmpattern{scale=1.5}{2}{1/2,2/1}{1/2}{1/0/2/1/1},
C)\decpatternww{scale=1.5}{2}{1/2,2/1}{}{}{1/0,1/2}{1/1/2/2/{\mpattern{scale=0.4}{2}{1/1,2/2}{}}}{}
\end{center}
Se invece ci si vuole assicurare che una coppia disordinata venga ordinata, la decorazione da apllicare consiste semplicemente nella "negazione" di tutte quelle espresse finora:
\begin{center}
\mmpattern{scale=1.7}{2}{1/2,2/1}{}{}$\Rightarrow$
\decpatternww{scale=1.7}{2}{1/2,2/1}{}{}{1/0,1/2}{}{1/1/2/2/{\mpattern{scale=0.4}{2}{1/1,2/2}{}}}
\end{center}
\subsection*{Composizione di {stacksort} con {POP-stacksort}}
$$(S{S_{POP}})^{-1}(Av(21)) = S_{POP}^{-1}(S^{-1}(Av(21))) = S_{POP}^{-1}(Av(231))$$
\\Si applicano le condizioni per le controimmagini ai due candidati: 231 e 321.\\\\
Nel caso di 231, le inversioni sono $(2,1),(3,1)$. Si osserva che 3 gi\'a soddisdfa la condzione A per la coppia $(2,1)$, quindi resta da applicare le condizioni alla coppia $(3,1)$:
\begin{center}
\mmpattern{scale=1.5}{3}{1/2,2/3,3/1}{}{2/3/3/4/1},
\mmpattern{scale=1.5}{3}{1/2,2/3,3/1}{2/3}{2/0/3/1/1},
\decpatternww{scale=1.5}{3}{1/2,2/3,3/1}{}{}{2/0,2/3}{2/1/3/3/{\mpattern{scale=0.4}{2}{1/1,2/2}{}}}{}
\end{center}
Dalla condizione A si ottiene il pattern 2341, dalla B il 3412 dalla condizione C si ottengono, considerando la posizione relativa degli elemeneti nella coppia aggiuntiva rispetto all'elemento 2, i pattern 25341, 35241, 45231.\\\\
Se gli elementi del pattern 231 nell'immagine invece formano un pattern 321 nella controimmagine, allora si applica la condizione per ordinare una coppia alla coppia $(3,2)$:
\begin{center}
\decpatternww{scale=2}{3}{1/3,2/2,3/1}{}{}{1/0,1/1,1/3}{}{1/2/2/3/{\mpattern{scale=0.5}{2}{1/1,2/2}{}}}
\end{center}
Si applicano adesso le condizioni per la coppia $(3,1)$ che ci si aspetta resti disordinata:
\begin{center}
\decpatternww{scale=2}{3}{1/3,2/2,3/1}{}{}{1/0,1/1,1/3}{2/3/3/4/1}{1/2/2/3/{\mpattern{scale=0.5}{2}{1/1,2/2}{}}}
\decpatternww{scale=2}{3}{1/3,2/2,3/1}{}{}{1/0,1/1,1/3,2/3}{2/0/3/1}{1/2/2/3/{\mpattern{scale=0.5}{2}{1/1,2/2}{}}}
\decpatternww{scale=2}{3}{1/3,2/2,3/1}{}{}{1/0,1/1,1/3,2/3,2/0}{2/1/3/3/{\mpattern{scale=0.5}{2}{1/1,2/2}{}}}{1/2/2/3/{\mpattern{scale=0.5}{2}{1/1,2/2}{}}}
\end{center}
Si osservi come applicando le condzioni A e B per la coppia $(3,1)$ queste siano gi\'a soddisfatte per la coppia $(2,1)$. Applicando la condizione C la coppia pu\'o essere costruita utilizzando l'elemento 2 come membro minore all'interno della coppia, e questo soddisfa anche la condizione A per la coppia $(2,1)$. Si ottengono dunque i pattern 3241, 4312, 4231.\\\\
Si osserva che, tra i risultati ottenuti prima, alcuni risultano non necessari: $2341\preceq25341, 3241\preceq35241, 4231\preceq45231$. Si semplificano e si ottiene il seguente risultato:
$$(S{S_{POP}})^{-1}(Av(21)) = Av(2341, 3241, 3412, 4231, 4312)$$
\subsection*{Composizione di {queuesort} con {POP-stacksort}}
$$Q{S_{POP}}^{-1}(Av(21)) = S_{POP}^{-1}(Q^{-1}(Av(21))) = S_{POP}^{-1}(Av(321))$$
\\Si applicano anche stavolta le condizioni per la ricerca di controimmagini al pattern 321. Essendo 321 una sequenza decrementale, non ci sono altri candidati; le sue inversioni sono $(3,2),(3,1),(2,1)$ ed \'e necessario applicare tutte le combinazioni possibili di condizioni A,B,C a queste coppie.\\\\
Una volta applicate le diverse condizioni alla coppia $(3,2)$ \'e evidente che queste soddisfano anche le condizioni per la coppia $(3,1)$:
\begin{description}
	\item applicando la condizione A e ponendo quindi un elemento maggiore di 3 tra 3 e 2, allora lo stesso elemento rappresenta un valore maggiore di 3 posizionato tra 3 e 1 (condzione A per la coppia $(3,1)$)
	\item se si applica la condizione B, ponendo un valore $2^-<2$ tra 2 e 3, allora o si ha che $2^-<1$, che soddisfa la condizione B per la coppia $(3,1)$, o che $1<2^-<2$, e allora la coppia $(2^-,2)$ \'e una coppia ordinata tra 3 e 1 (condizione C per ls coppia $(3,1)$)
	\item nel caso, infine, in cui si applichi la condizione C e si ponga quindi una coppia ordinata tra 3 e 2 allora la stessa coppia soddisfa la condizione C per la coppia $(3,1)$.
\end{description}
In sintesi, i seguenti pattern soddisfano le condzioni da imporre alle coppie $(3,2),(3,1)$:
\begin{center}
\mmpattern{scale=1.7}{3}{1/3,2/2,3/1}{}{1/3/2/4/1}
\mmpattern{scale=1.7}{3}{1/3,2/2,3/1}{1/3}{1/0/2/2/1}
\decpatternww{scale=1.7}{3}{1/3,2/2,3/1}{}{}{2/0,2/2,2/3}{2/1/3/2/{\mpattern{scale=0.5}{2}{1/1,2/2}{}}}{}
\end{center}
Una volta applicate le condizioni anche alla coppia $(2,1)$ si ottengono i seguenti 9 meta-pattern, da cui si ottengono 22 pattern classici:
\begin{center}
\mmpattern{scale=1.7}{3}{1/3,2/2,3/1}{}{1/3/2/4/1,2/2/3/4/1}
\mmpattern{scale=1.7}{3}{1/3,2/2,3/1}{2/3,2/2}{1/3/2/4/1,2/0/3/1/1}
\decpatternww{scale=1.7}{3}{1/3,2/2,3/1}{1/3/2/4/1}{}{2/0,2/2,2/3}{2/1/3/2/{\mpattern{scale=0.5}{2}{1/1,2/2}{}}}{}
\mmpattern{scale=1.7}{3}{1/3,2/2,3/1}{1/3}{1/0/2/2/1,2/2/3/4/1}
\mmpattern{scale=1.7}{3}{1/3,2/2,3/1}{2/2,2/3}{1/0/2/2/1,2/0/3/1/1}
\decpatternww{scale=1.7}{3}{1/3,2/2,3/1}{}{}{1/3,2/3,2/2,2/0}{1/0/2/2/1,2/1/3/2/{\mpattern{scale=0.5}{2}{1/1,2/2}{}}}{}
\decpatternww{scale=1.7}{3}{1/3,2/2,3/1}{}{}{1/0,1/1,1/3}{2/2/3/4/1,1/2/2/3/{\mpattern{scale=0.5}{2}{1/1,2/2}{}}}{}
\decpatternww{scale=1.7}{3}{1/3,2/2,3/1}{}{}{1/0,1/1,1/3,2/2,2/3}{2/0/3/1/1,1/2/2/3/{\mpattern{scale=0.5}{2}{1/1,2/2}{}}}{}
\decpatternww{scale=1.7}{3}{1/3,2/2,3/1}{}{}{1/0,1/1,1/3,2/0,2/2,2/3}{1/2/2/3/{\mpattern{scale=0.5}{2}{1/1,2/2}{}},2/1/3/2/{\mpattern{scale=0.5}{2}{1/1,2/2}{}}}{}
\end{center}
Scrivendo come pattern classici quelli che si ottengono da questi meta-patterns ed eliminando quelli ridondanti, si ottiene il seguente risultato:
\begin{eqnarray*}
&Q{S_{POP}}(Av(21)) = Av(34251, 35241, 41352, 42351, 45231, \\
& 45312, 51342, 51423, 52341, 52413, 53412)
\end{eqnarray*}
\subsection*{Composizione di {bubblesort} con {POP-stacksort}}
$$(B\circ{S}_{POP})^{-1}(Av(21)) = S_{POP}^{-1}(B^{-1}(Av(21))) = S_{POP}^{-1}(Av(231,321))$$
Il l'insieme minimo di pattern per l'ordinabilit\'a di permutazione con la composzione bubblesort, POPstacksort \'e l'unione degli insiemi dei due risultati precedenti, stacksort e queuesort.\\
Si nota, tuttavia che ogni pattern contenuto nell'insieme ottenuto dalla combinazione con queuesort contiene \textbf{almeno} un pattern tra quelli ottenuti da stacksort. Nello specifico, a titolo esemplificativo:$$3241\preceq34251, 35241, 41352, 42351$$$$3412\preceq53412$$$$4231\preceq45231, 51342, 52341, 52413$$$$4312\preceq45312, 51423$$ 
Dunque l'insieme minimo di pattern che caratterizza le permutazioni ordinabili da $B\circ{S}_{POP}$ \'e lo stesso di $S\circ{S}_{POP}$.
$$(B{S_{POP}})^{-1}(Av(21)) = Av(2341, 3241, 3412, 4231, 4312)$$
\subsection*{Composizione di stacksort con POP-queuesort}
[ancora da correggere]
Grazie a vari risultati ottenuti applicando i software mostrati nel capitolo 3, \'e stato possibile formulare la seguente congettura:
$$Cons^{-1}Av(231) = Av(2431, 4231, 23514, \mmpattern{scale=0.7}{5}{1/3,2/2,3/5,4/1,5/4}{0/4,0/3,0/5,1/3,1/4,1/5}{})$$
Per dimostrare la congettura occorre dimostrare la doppia inclusione:
\begin{description}
\item[(A)]$Cons^{-1}Av(231) \subseteq Av(2431, 4231, 23514, \mmpattern{scale=0.7}{5}{1/3,2/2,3/5,4/1,5/4}{0/4,0/3,0/5,1/3,1/4,1/5}{})$
\item[(B)]$Av(2431, 4231, 23514, \mmpattern{scale=0.7}{5}{1/3,2/2,3/5,4/1,5/4}{0/4,0/3,0/5,1/3,1/4,1/5}{}) \subseteq Cons^{-1}Av(231)$
\end{description}
\paragraph*{Nota:}Per tutta la dimostrazione i valori che formano i pattern saranno indicati da $a,b,c,d$ ed $e$ (quando presente), $n$ indicher\'a la lunghezza della permutazione $\pi$ e si intende $1\leq{a}<b<c<d<e\leq{n}$.
\paragraph*{Osservazione:} Sia $\pi$ una permutazione in input all'algoritmo \textit{Cons}, tutti gli elementi di $\pi$ che non sono LTR massimi effettuano un bypass. Un LTR massimo pu\'o accodarsi se e solo se \'e il consecutivo dell'elemento in fondo alla coda e provoca un \textit{pop} in tutti gli altri casi.
\begin{center}
\textbf{(A)} $Cons^{-1}Av(231) \subseteq Av(2431, 4231, 23514, \mmpattern{scale=0.7}{5}{1/3,2/2,3/5,4/1,5/4}{0/4,0/3,0/5,1/3,1/4,1/5}{})$
\end{center}
\begin{proof}
Dobbiamo dimostrare che tutte le permutazioni la cui immagine evita il pattern 231 evitano i pattern $2431, 4231, 23514, \mmpattern{scale=0.7}{5}{1/3,2/2,3/5,4/1,5/4}{0/4,0/3,0/5,1/3,1/4,1/5}{}$, ovvero che l'immagine di ogni permutazione che contiene almeno uno di questi pattern contiene un pattern 231.
\paragraph*{2431}
$$\pi = \dots{b}\dots{d}\dots{c}\dots{a}\dots$$
Le non inversioni $(b,d)$ e $(b,c)$ saranno presenti anche nell'immagine e inoltre $c,a$, non essendo LTR massimi, effettuano un bypass. Si osserva che se $b$ effettua un bypass si ha necessariamente un pattern 231 nella sottosequnze $bca$. Se $b$ entra in coda allora $d$ provoca necessariamentee un $pop$ (non pu\'o essere accodato alla stessa coda che contiene $b$, perch\'e la coda \'e consecutiva e $c$, che \'e un valore intermedio tra $b$ e $d$, \'e pi\'u avanti nell'input) e l'immagine contiene la sottoseqenza $bca$.
\paragraph*{4231}
$$\pi = \dots{d}\dots{b}\dots{c}\dots{a}\dots$$
Nessuno degli elementi che formano la sottosequenza $bca$ contenuta in $\pi$ \'e un LTR massimo, quindi bypassano tutti e la sequanza \'e contenuta anche in $Cons(\pi)$.
\paragraph*{23514}$$\pi = \dots{b}\dots{c}\dots{e}\dots{a}\dots{d}\dots$$
Le non inversioni sono presenti anche nell'immagine, in particolare la non inversione $(b,c)$, inoltre $a$ non \'e un LTR massimo quindi effetua un bypass. Si vuole mostrare che la non inversione $(b,c)$ viene sempre inserita nell'output prima di $a$. Se $b,c$ effettuano un bypass si ha sicuramente $bca$ nell'immagine. Se entrambi entrano in coda almeno un pop viene provocato da $e$ prima che si arrivi ad $a$. Anche se $b$ effettua un bypass e $c$ entra in coda quest'ultimo viene sicuramente aggiunto all'output prima di $a$, eventualmente da $e$.   
\subsection*{$\mmpattern{scale=1}{5}{1/3,2/2,3/5,4/1,5/4}{0/4,0/3,0/5,1/3,1/4,1/5}{}$}
Una permutazione che contiene occorrenza di questo pattern ha la forma:
$$\pi = \dots{c}\dots{b}\dots{e}\dots{a}\dots{d}\dots$$dove nessun valore precedente a $b$ \'e maggiore di $c$.\\\\
Si osserva che $b,a$ effettuano un bypass. $c$ \'e un LTR massimo, quindi sia che provochi un pop o meno, verr\'a inserito in coda, dopodich\'e si verifica il bypass di $b$, e $c$ viene sicuramente estratto dalla coda prima di arrivare ad $a$ (eventualmente \'e $e$ a provocare il pop). Si pu\'o escludere che $c$ esca dalla coda prima che $b$ effettui un bypass perch\'e prima di $b$ nessun elemento \'e maggiore di $c$.\\Si forma cos\'i la sottosequenza $bca$.
\end{proof}
\begin{center}
\textbf{(B)} $Av(2431, 4231, 23514, \mmpattern{scale=0.7}{5}{1/3,2/2,3/5,4/1,5/4}{0/4,0/3,0/5,1/3,1/4,1/5}{}) \subseteq Cons^{-1}Av(231)$
\end{center}
\begin{proof}
Dobbiamo dimostrare che tutte le permutazioni che evitano i pattern $2431, 4231, 23514, \mmpattern{scale=0.7}{5}{1/3,2/2,3/5,4/1,5/4}{0/4,0/3,0/5,1/3,1/4,1/5}{}$ sono contenute anche nell'insieme $Cons^{-1}Av(231)$, ovvero si dimostrer\'a che ogni permutazione $\pi$ tale che $231\preceq{Cons}(\pi)$ contiene uno dei pattern elencati.\\
Si osserva che se 3 elementi nell'immagine $Cons(\pi)$ formano un pattern 231, quegli stessi elementi devono formare un pattern 231 o 321 nella preimmagine, dato che \textit{Cons} non produce nuove inversioni. Si analizzano questi due casi separatamente.
\paragraph*{$231\preceq\pi$:}
$$\pi=\dots b\dots c\dots a\dots$$
$$Cons(\pi)=\dots b\dots c\dots a\dots$$
La non inversione $(b,c)$ resta anche nell'imamgine. Si vuole definire l'ordine relativo degli altri elementi di $\pi$ rispetto a $a,b,c$ perch\'e le inversioni $(b,a),(c,a)$ non vengano rimosse da \textit{Cons}.\\\\
Se la forma di $\pi$ \'e tale che $b$ entra in coda, per mentenere la sottosequenza $bca$ \'e necessario che $c$ non sia un LTR massimo. Questa condizione \'e soddisfatta se $\pi$ contiene un'occorrenza del mesh-pattern \mmpattern{scale=0.7}{3}{1/2,2/3,3/1}{}{0/3/2/4/1}, ovvero se contiene uno dei pattern classici 4231 o 2431.\\\\
Se la forma di $\pi$ \'e invece tale per cui $c$ entri in coda, devono essere presenti degli elementi in $\pi$ che provocano un pop prima che $a$ possa bypassare (si noti come queste considerazioni si applicano sia nel caso in cui $b$ sia in coda con $c$ al momento del pop, sia in quelli in cui $b$ \'e gi\'a stato inserito nell'output). Per assicurare il pop \'e sono necessari due nuovi elementi: $e,d$ dove $e$ deve essere posizionato tra $c$ e $a$ mentre $d$ dopo di $e$. L'elemento $e$ assicura il pop se \'e maggiore della testa della coda (se \'e minore tutti gli elementi del pattern effettuano un bypass) ma non pu\'o accodarsi, in quanto $d$ non pu\'o ancora essere in coda e quindi $e$ non pu\'o essere consecutivo alla coda. Queste condizioni sono soddisfatte dal mesh-pattern $\mmpattern{scale=0.7}{4}{1/2,2/3,3/4,4/1}{}{3/3/5/4/1}$, quindi si deve avere che $\pi$ contenga almeno uno tra i pattern classici 23541,23514. Il primo dei due pattern classici viene ignorato, in quanto contiene uno dei pattern gi\'a trovati prima:$2431\preceq23541$.  
\paragraph*{$321\preceq\pi$:}$$\pi=\dots c\dots b\dots a\dots$$$$Cons(\pi)=\dots b\dots c\dots a\dots$$
Ci chiediamo quale deve essere l'ordine relativo degli altri elementi di $\pi$ rispetto a $a,b,c$ per ottenere che l'inversione $(3,2)$ venga ordinata mentre le altre restino invariate.\\\\
Per ordinare $(b,c)$, $c$ deve essere un LTR massimo, in modo che possa entrare nella coda, e non deve verificarsi nessun pop tra $c$ e $b$. Queste condizioni sono rappresentate dalle aree oscurate nel mesh pattern \mmpattern{scale=0.7}{3}{1/3,2/2,3/1}{0/3,1/3}{}. Per forzare un pop tra $c$ e $a$ \'e necessario introdurre due elementi $d,e$ come nel caso precedente: si ottiene cos\'i il mesh-pattern \mmpattern{scale=0.7}{4}{1/3,2/2,3/4,4/1}{0/3,1/3,0/4,1/4}{3/3/5/4/1}, la cui parte decorata permette rappresenta l'unione delle possibili posizioni del valore $d$ nei due successivi mesh-pattern mostrati:
\begin{center}
\mmpattern{scale=1.5}{4}{1/3,2/2,3/4,4/1}{0/3,1/3,0/4,1/4}{3/3/5/4/1} $\Rightarrow$
\mmpattern{scale=1.2}{5}{1/3,2/2,3/5,4/1,5/4}{0/3,1/3,0/4,1/4,0/5,1/5}{},
\mmpattern{scale=1.2}{5}{1/3,2/2,3/5,4/4,5/1}{0/3,1/3,0/4,1/4,0/5,1/5}{} 
\end{center}
Degli ultimi due mesh pattern mostrati, si ignora il secondo in quanto contiene un'occorrenza del pattern 2431, gi\'a trovato in precedenza, nella sottosequenza $beda$.
\end{proof}
\subsection*{$Q\circ{Q_{POP}}$}
[Congettura per la controimmagine di 321, per ora scritta solo su carta]
\subsection*{$B\circ{Q_{POP}}$}
[Unire risultati dei due paragrafi precedenti]

\bibliographystyle{plain}
\bibliography{Chapters/bibliografia}
%--------------------------------------------------------------
\end{document}
%--------------------------------------------------------------
